\clearpage
\VerdierPhysicalPageAnchor{0228}%
\begingroup\small\setlength{\parskip}{0pt}\setlength{\emergencystretch}{1.5em}
\setlength{\abovedisplayskip}{4pt plus 2pt minus 2pt}%
\setlength{\belowdisplayskip}{4pt plus 2pt minus 2pt}%

\VerdierSemanticAnchor{ch4d-4-1-4}
\medskip\noindent\textbf{4.1.4.}\quad
(4.1.3) より、$\operatorname{comp}(\mathcal A)$ の任意の対象 $Y$ に対し、

\VerdierSemanticAnchor{ch4d-4-1-4-spec-k}
\[
\operatorname{Spec}_1(Y)=\bigl(Y(p,q),\delta(p,q,r)\bigr),
\]
は (4.1.2.3)、(4.1.2.4)、(4.1.2.6) で定義され、$K(\mathcal A)$ に値を
とるスペクトル対象である。スペクトル対象 $\operatorname{Spec}_1(Y)$ は
$\operatorname{comp}(\mathcal A)$ の対象 $Y$ に関手的に依存する。したがって
関手

\VerdierSemanticAnchor{ch4d-4-1-4-functor}
\[
Y\longmapsto\operatorname{Spec}_1(Y)
\]
が $\mathcal A$ の複体の圏から $K(\mathcal A)$ に値をとるスペクトル対象の
圏へ定まる。この関手を\emph{$K(\mathcal A)$ に値をとる第一標準スペクトル
対象関手}という。$\operatorname{comp}(\mathcal A)$ の各対象 $Y$ に対し、
スペクトル対象 $\operatorname{Spec}_1(Y)$ を\emph{複体 $Y$ に付随する、
$K(\mathcal A)$ に値をとる第一スペクトル対象}という。$Y$ が
$\operatorname{comp}^{*}(\mathcal A)$ の対象であり、$*$ が $+$、$-$、$b$、
または「空」のいずれかであるとき、複体 $Y$ に付随する、$K(\mathcal A)$
に値をとる第一スペクトル対象は、実際には $K^{*}(\mathcal A)$ に値をとる
ことに注意する。

\VerdierSemanticAnchor{ch4d-4-1-5}
\medskip\noindent\textbf{4.1.5.}\quad
$\mathcal A$ をアーベル圏とし、$Y$ を $\mathcal A$ の複体とする。複体
$Y$ に付随する第一スペクトル対象の標準関手
$K(\mathcal A)\longrightarrow D(\mathcal A)$ による像は、
$D(\mathcal A)$ に値をとるスペクトル対象であり、\emph{複体 $Y$ に付随する、
$D(\mathcal A)$ に値をとる第一スペクトル対象}と呼ばれる。多くの場合、
単に\emph{複体 $Y$ に付随する第一スペクトル対象}といい、記法を濫用して

\VerdierSemanticAnchor{ch4d-4-1-5-spec-d}
\[
\operatorname{Spec}_1(Y)=\bigl(Y(p,q),\delta(p,q,r)\bigr)
\]
と書く。これは複体 $Y$ に付随する、$D(\mathcal A)$ に値をとる第一
スペクトル対象を表す。第一標準スペクトル対象関手と標準関手
$K(\mathcal A)\longrightarrow D(\mathcal A)$（スペクトル対象へ拡張したもの）
との合成を、\emph{$D(\mathcal A)$ に値をとる第一標準スペクトル対象関手}
という。これは $\mathcal A$ の複体の圏から $D(\mathcal A)$ のスペクトル
対象の圏への関手であり、\emph{一般には、この関手は複体のホモトピー圏を
経由して分解しない}ことを強調しておく。

\VerdierSemanticAnchor{ch4d-4-1-6}
\medskip\noindent\textbf{4.1.6.}\quad
$\mathcal A$ の対象の $n$-重複複体から、フィルトレーションを備えた単純複体、
したがって $K(\mathcal A)$ に値をとるスペクトル対象を構成できる。$X$ を
$n$-重複複体（第 I 章、2.1.2）とし、$m$ を整数、$1\leq m\leq n$ とする。
各 $\sigma\in\mathbb Z[n]$（第 I 章、2.1.1 の記法）について、その第
$m$ 座標を $\sigma_m$ と書く。

\VerdierSemanticAnchor{ch4d-4-1-6-family-setup}
まず、複体と複体の射の族を定義する。
$\mathcal{F}\ell_1(\widetilde{\mathbb Z})$ の各組 $(p,q)$ に対し、
${}^{m}X(p,q)$ を
\endgroup
